\documentclass[a4paper,twoside,12pt]{report}
\usepackage[utf8x]{inputenc}
\usepackage{color}
\usepackage{fancyvrb} % pour mettre les verbatim dans des boites
%\usepackage[dvipdf]{graphics}
\usepackage[pdftex]{graphicx}
\usepackage{fancyhdr}
\usepackage{listings}
\usepackage{float}
\usepackage{listingsutf8}
\usepackage[T1]{fontenc}
\usepackage[french]{babel}
\usepackage{textcomp}
\usepackage[toc,page]{appendix}
%\usepackage{hyperref}
% Quelques réglages particuliers
\oddsidemargin=-0.1in %xx
\evensidemargin=-0.1in
\textwidth=6.1in
\topmargin=-0.5in
\textheight=8.7in
\parskip 0.25in
\lstset{
language=Java,
basicstyle=\normalsize,
upquote=true,
aboveskip={1.5\baselineskip},
columns=fullflexible,
showstringspaces=false,
extendedchars=true,
breaklines=true,
showtabs=false,
showspaces=false,
showstringspaces=false,
identifierstyle=\ttfamily,
keywordstyle=\color[rgb]{1,0,1},
commentstyle=\color[rgb]{0.133,0.545,0.133},
stringstyle=\color[rgb]{0.627,0.126,0.941},
}
\fancyhead[RO,RE]{\includegraphics[width=0.5in]{Logo-univ-orleans.png}}
\fancyhead[CO,CE]{Rapport de projet}
\fancyhead[LO,LE]{Jeu de Taquin}
\fancyfoot[C]{2014}
\fancyfoot[RO, LE] {\thepage}
%\fancyfoot[LO, RE] { ARNOULT S., MEKHILEF W., OUSSAD J., RETY M.}
%
%Quelques macros
\newcommand{\moncode}[1]{\begin{center}
                        \lstinputlisting[inputencoding=utf8/latin1]{#1}
                        \end{center}}

\newcommand{\monimage}[3]{
\par\noindent
\begin{figure}[H] %on ouvre l'environnement figure
\begin{center}
\includegraphics[width=0.7\textwidth]{#1} %ou image.png, .jpeg etc.
\caption{#2} %la légende
\label{#3} %l'étiquette pour faire référence à cette image
\end{center}
\end{figure} %on ferme l'environnement figure
}

\newcommand{\ml}[0]{\par\noindent}

%\hyphenation{supplémentaires}{sup-pl\'e-men-tai-res}
%opening
\title{\textcolor{blue}{\Large Rapport du projet IF06}\\\textcolor{blue}{\Large Jeu de Taquin}}
\vskip 1cm
\author{ARNOULT Simon, MEKHILEF Wissame, OUSSAD Jihad, RETY Martin \date{\today}}


\begin{document}

\thispagestyle{empty}
%
\begin{figure}[H]
\includegraphics[width=0.2\linewidth]{Logo-univ-orleans.png}
% \hfill

\end{figure}
\vspace{2cm}
%
\begin{center}
{\Huge Licence 2 Informatique\\\ \\Rapport du projet IF04}
\par\vspace{1.4cm}

{\Huge\bf \textcolor{red}{\bf Jeu de Taquin}}
\par\vspace{1.6cm}

{\Large       Réalisé par:}
\par\vspace{1.3cm}
{\large\bf \textcolor{blue}{ARNOULT Simon, MEKHILEF Wissame, OUSSAD Jihad, RETY Martin}}
\vfill
\today
\end{center}
\newpage
\pagestyle{fancy}

\begin{abstract}
%
Nous avons le plaisir de vous présenter notre travail sur ce projet. Durant la deuxième année de la l2 Informatique
à l'Université d'Orléans nous avons travaillé en groupe de quatre sur un projet de résolution de Taquin.
\end{abstract}
 
\newpage
\tableofcontents
\listoffigures
\newpage

\chapter{Introduction}
Ce rapport s'articulera autour de quatre chapitre, nous allons détailler notre travail dans l'ordre chronologique. En commençant par l'étude
du projet, puis nous allons détailler notre organisation pour le travail de groupe, nous continuerons ensuite sur le développement et nous
finirons ce rapport par une analyse des algorithmes.
%
\section{Cahier des charges}

%
\chapter{Etude}
L'étude du projet fut une première étape importante pour correctement se mettre dans une bonne dynamique de groupe. Dés la première semaine
nous avons réaliser une version jouable du Taquin, nous avons travaillé sur cette version durant toute la première phase d'essaie.
Rapidement il a fallu faire des essaies sur ce jeu pour qu'il puisse se résoudre algorithmiquement. Nous avons chacun travailler de notre
coté sur nos idées pour pouvoir en tester un plus grand nombre, mais tout en restant en contact réguliérement pour avancer un ensemble.
Cette manière de répartir les tâches nous à permi de se rendre compte des difficulté qu'on allé rencontrer plus rapidement que si on avais
tout fais ensemble.
%
\section{Analyse de faisabilité}
L'analyse de faisabilité fut un moment qui à durer du début du projet jusqu'au vacances de février, durant cette phase chacun travaillant
sur des fonctionnalité différents, on pu voir ou était les problèmes dans notre architecture de départ. Ces tests on été trés diverses. On
a travailler sur les fonctionnalitées VT100 du terminal et la récupération des touches tappé par l'utilisateur. Mais on a aussi crée des
algorithmes pour essayer de résoudre le Taquin. Tout ces tests on permis au mois de février d'avoir une idée clair de l'architecture du
projet.
\section{Conception UML}
Cette analyse de février a permis d'aboutir au diagramme UML suivant, diagramme qui n'a pas beaucoup changé jusqu'à la version finale.
%INCLURE LE MODELE UML
\chapter{La gestion de projet}
%
La gestion de projet est à la base de tout projet encore plus quand il se fais avec une équipe de quatre personnes.
\section{Le travail de groupe}
%
La gestion de projet fut au coeur de nos préocupations avant même que le projet ne démarre, on a cherché à créer une équipe dynamique. On
c'est donc vu réguliérement dans les salles de l'Université. Mais
malgré celà il a fallut mettre en place des moyens dédiée pour faciliter le travail et éviter une dégradation de l'entente. Dès la
première semaine, on mis en place un dépôt git sur le site de l'hébergeur GitHub, ce dépôt et le faite de l'avoir dés la premiere heure
nous à permi d'avoir le réflexe de l'utiliser même si on rencontrer quelque problèmes, il a permis a chacun de voir l'avancement du projet.
Cependant un tel dépôt ne répond pas à la question de la communication au seins de l'équipe, chacun habitant une ville différente et Skype
nous à aidé dans cette tâche.
\section{Le diagramme de Gantt}
%
Nous avons utilisé la technique du diagramme de GANTT pour vérifier en interne que le projet avancé convenablement. Ci-dessous vous pouvez voir ce
diagramme.
\monimage{DiagrammeDeGantt.pdf}{Diagramme de GANTT}{DG}
\chapter{Phase de développement}
%
La phase de développement...
\section{L'architecture}
On explique les choix qu'on a fais
\section{Le codage}
%
\section{Junit et Benchmark}
%
\subsection{Utilisation des benchmarks}
On a utilisé les benchmark pour améliorer l'efficacité de l'algorithme de solution, en effet on a remarquer entre autre que la taille de l'ensemble
incomplet avais une importance dans le temps de résolution. On a donc créer un test avec plusieurs taille d'ensemble pour voir laquelle 
permetté une resolution rapide.
\monimage{ensembleincomplettest2.pdf}{Exemple}{EX}
\chapter{Analyse et Conclusion}
%
\section{Analyse des Benchmark}
%
\subsection{Jeu simple}
%
\subsection{Jeu moyen}
%
\subsection{Jeu compliqué}
%
\section{Conclusion}
%
\newpage
\bibliography{bibliographie}{}
\bibliographystyle{plain}

\chapter{Resources utilisées}
\begin{itemize}
 \item JUNIT
 \item JAVA
 \item Dia
 \item Eclipse
 \item Argouml
 \item \LaTeX (sous kile)
 \item Google Drive
 \item Kate
 \item Git, GitHub
\end{itemize}

\end{document}
